\documentclass[11pt,a4paper]{article}
\usepackage[utf8]{inputenc}
\usepackage[T1]{fontenc}
\usepackage[margin=2.5cm]{geometry}
\usepackage{amsmath,amssymb}
\usepackage{booktabs}
\usepackage{array}
\usepackage{hyperref}

\newcommand{\tx}[1]{\texttt{#1}}

\title{NSUCRYPTO 2026, Round 2\\Problem 1: ``An ancient civilization''}
\date{}
\author{Jacob Lu, Raymond Ouyang, Ruiyang Ye}
\begin{document}
\maketitle

\section*{Answer}

\begin{itemize}
  \item \textbf{Key:} \tx{SATOSHINAKAMOTO}
  \item \textbf{Cipher:} the 102 letters of Line~3, read across the seven tablets as one stream, form a repeating-key Vigen\`ere cipher that preserves case.
  \item \textbf{Message:}
\end{itemize}

\begin{quote}
The hawk was having a breakfast of cookies and sea buckthorn tea but was
caught by artificial intelligence and put in a cage.
\end{quote}

\section{The archive is a blockchain}

Each tablet has four lines. We read them as the fields of a block in a
Bitcoin-style chain:

\begin{center}
\begin{tabular}{@{}lll@{}}
\toprule
Line & Content & Role \\
\midrule
1 & $0,1,\dots,6$ & block index \\
2 & 64 zeros on tablet~0, empty elsewhere & previous-block hash \\
3 & 16 letters, with 6 on the last tablet & block data \\
4 & a decimal integer & nonce \\
\bottomrule
\end{tabular}
\end{center}

Several features support this reading:

\begin{itemize}
  \item A string of 64 hexadecimal zeros is the
  previous-block hash of Bitcoin's genesis block, written as 64 hex characters.
  Only block~0 can carry it.
  \item ``The past is the key to the present. The
  present is the key to the future'' is conceptually analogous with hash chaining: each block
  commits to the one before it.
  \item The seven
  sectors are the seven blocks, and the 64 stars are the 64 hex characters of
  a hash.
  \item A proof-of-work search that
  needs $b$ leading zero bits takes a geometrically distributed number of
  attempts with mean~$2^b$. The seven numbers fit this exactly:
  \begin{itemize}
    \item The mean is $16\,617\,838 \approx 2^{23.99}$;
    \item The maximum-likelihood difficulty under the geometric model is
    $23.99$ bits (95\% CI $23.1$--$25.3$);
    \item A Kolmogorov--Smirnov test against an exponential distribution with
    mean $2^{24}$ gives $p=0.33$, while means $2^{20}$ and $2^{28}$ are
    rejected ($p<0.001$);
    \item The geometric model beats a uniform one by a likelihood ratio of
    about $14:1$.
  \end{itemize}
  Because target of 24 zero bits is four zero characters in base~64, which made this a clue corresponding to ``the first four galaxies''.
\end{itemize}

The archive is therefore a model of the Bitcoin blockchain. The ``key'' that
the inscription points to is the name of its creator, \tx{SATOSHINAKAMOTO}.

\section{Decryption}

The tablets are shown side by side, with the first \tx{L} being capitalized. So Line~3 is read straight across as a
single stream of $6\cdot 16+6=102$ letters:
\begin{center}
\tx{Lhxvsdsjachmjbby | aufwhssactatvcgk | bskhvqsoanivylhh | ffamnbetiolqsuzv}\\
\tx{lignrdirwvwslbbl | ltyiqezqxofdiilp | vnckgq}
\end{center}

Number the letters $A=0,\dots,Z=25$, and let $k_0k_1\dots k_{14}$ be the key
\tx{SATOSHINAKAMOTO}. Each ciphertext letter $c_i$ decrypts to
\[
p_i \equiv c_i - k_{i \bmod 15} \pmod{26}, \qquad i=0,1,\dots,101,
\]
with case kept: the capital \tx{L} at the start becomes the capital \tx{T}.

The key has length 15, while the tablets are 16 letters wide. It therefore
runs continuously across the tablets and does not restart at each one.
Grouped into key periods:

\begin{center}
\renewcommand{\arraystretch}{1.15}
\begin{tabular}{@{}r>{\ttfamily}l>{\ttfamily}l>{\ttfamily}l@{}}
\toprule
Letters & \normalfont Ciphertext & \normalfont Key & \normalfont Plaintext \\
\midrule
1--15   & Lhxvsdsjachmjbb & SATOSHINAKAMOTO & Thehawkwashavin \\
16--30  & yaufwhssactatvc & SATOSHINAKAMOTO & gabreakfastofco \\
31--45  & gkbskhvqsoanivy & SATOSHINAKAMOTO & okiesandseabuck \\
46--60  & lhhffamnbetiolq & SATOSHINAKAMOTO & thornteabutwasc \\
61--75  & suzvlignrdirwvw & SATOSHINAKAMOTO & aughtbyartifici \\
76--90  & slbblltyiqezqxo & SATOSHINAKAMOTO & alintelligencea \\
91--102 & fdiilpvnckgq    & SATOSHINAKAMO   & ndputinacage \\
\bottomrule
\end{tabular}
\end{center}

\noindent
Example: $c_0=\tx{L}=11$ and $k_0=\tx{S}=18$, so $p_0\equiv 11-18\equiv 19
\pmod{26}$, which is \tx{T}.


\noindent
With spaces added it reads:
The hawk was having a breakfast of cookies and sea buckthorn tea but was
caught by artificial intelligence and put in a cage.

All 102 letters decrypt to grammatical English, which confirms the key.

\end{document}
